% evaluation.tex — Section 5: Evaluation
% Included by main.tex as % evaluation.tex — Section 5: Evaluation
% Included by main.tex as % evaluation.tex — Section 5: Evaluation
% Included by main.tex as % evaluation.tex — Section 5: Evaluation
% Included by main.tex as \input{sections/evaluation}
% No \begin{document} / \end{document}

\subsection{Methodology}
\label{sec:eval:method}

\paragraph{Hardware.}
All benchmarks were conducted on a workstation equipped with an NVIDIA RTX A6000
GPU (48~GB GDDR6 VRAM, Ampere architecture, 309~GB/s memory bandwidth).
The host runs Ubuntu 22.04 with CUDA 12.2 and driver 535.

\paragraph{Software stack.}
The VisionServe Go HTTP server was built with Go~1.22 and linked against
ONNX Runtime~1.17 \cite{onnxruntime} with the CUDA Execution Provider enabled.
Model artefacts are exported ONNX files; no TensorRT engine cache was used so
that results are reproducible across machines without a TensorRT installation.
The ONNX format and runtime are described in \cite{onnx2019}.

\paragraph{Benchmark protocol.}
For each model, five warm-up requests were issued and discarded to ensure the
ONNX session was fully initialised and any JIT optimisation passes had completed.
Twenty timed requests were then issued sequentially; all measurements are reported
over these 20 requests.  Requests were sent over the loopback interface
(\texttt{localhost:11435}) using a single-threaded HTTP client to eliminate
network variance.

\paragraph{Metrics.}
\begin{itemize}
  \item \textbf{p50 / p95} — 50th and 95th percentile of \emph{end-to-end HTTP
    latency} (wall-clock time from client send to response received), in milliseconds.
    This includes HTTP serialisation, Go image decoding and preprocessing, ONNX
    inference, postprocessing, and JSON serialisation.
  \item \textbf{srv\_p50} — 50th percentile of \emph{server-side inference latency},
    defined as the time spent exclusively inside ONNX session \texttt{Run} calls.
    This isolates pure tensor-computation cost from Go-layer overhead.
  \item \textbf{RPS} — sustained throughput in requests per second at batch size~1,
    derived from the mean end-to-end latency.
  \item \textbf{VRAM (MB)} — peak GPU memory as reported by \texttt{nvidia-smi
    --query-compute-apps} at inference time, after the session is warm.
  \item \textbf{ONNX (MB)} — size of the ONNX model file(s) on disk.
  \item \textbf{Cold-start (s)} — elapsed time from process start (before any request)
    to the completion of the first successful inference, measuring model load from disk,
    ONNX session initialisation, and provider graph optimisation.
\end{itemize}

\subsection{Latency and Throughput}
\label{sec:eval:latency}

Table~\ref{tab:benchmark} reports the full benchmark results for all 12 models with
available measurements, sorted by p50 latency ascending.  A horizontal rule separates
lightweight models (p50 $<$ 100~ms) from heavier models (p50 $\geq$ 100~ms).

\begin{table*}[t]
  \centering
  \small
  \caption{%
    VisionServe benchmark results on NVIDIA RTX A6000 (48~GB VRAM).
    20 timed requests after 5 warm-up discards.
    Sorted by p50 ascending; horizontal rule separates lightweight ($<$100~ms)
    from heavy ($\geq$100~ms) models.
    Four models (rt-detr, depth-anything-v2, nano-sam, grounded-sam) are
    not yet benchmarked and are omitted.%
  }
  \label{tab:benchmark}
  \begin{tabular}{llrrrrrrr}
    \toprule
    \textbf{Model} & \textbf{Task}
      & \textbf{p50 (ms)} & \textbf{p95 (ms)}
      & \textbf{srv\_p50 (ms)} & \textbf{RPS}
      & \textbf{VRAM (MB)} & \textbf{ONNX (MB)} & \textbf{Cold (s)} \\
    \midrule
    \texttt{clip}          & embed          &  33 &  69 &  12 & 27.9 &  810 & 335 & 5.8 \\
    \texttt{mobilenet-v3}  & classification &  38 &  71 &  18 & 25.4 &  308 &  16 & 2.1 \\
    \texttt{efficientnet-b0}& classification&  40 &  75 &  21 & 23.8 &  356 &  21 & 2.4 \\
    \texttt{scrfd}         & detection      &  45 &  69 &  23 & 22.4 &  420 &  16 & 3.9 \\
    \texttt{paddle-ocr}    & detection      &  54 &  78 &  34 & 17.3 &  462 &  15 & 4.7 \\
    \texttt{rf-detr-nano}  & detection      &  57 &  89 &  37 & 16.5 &  548 &  62 & 4.2 \\
    \texttt{midas}         & depth          &  65 &  98 &  44 & 14.7 &  420 &  48 & 3.8 \\
    \texttt{rf-detr}       & detection      &  78 & 112 &  56 & 12.3 &  804 & 168 & 5.1 \\
    \midrule
    \texttt{mobile-sam}    & segmentation   & 161 & 224 & 138 &  6.0 &  966 &  38 & 4.9 \\
    \texttt{efficient-sam} & segmentation   & 181 & 247 & 158 &  5.5 & 1628 &  39 & 5.8 \\
    \texttt{sam2}          & segmentation   & 242 & 544 & 222 &  3.9 & 2508 & 148 & 5.3 \\
    \texttt{grounding-dino}& open\_vocab    & 570 & 812 & 548 &  1.7 & 4392 & 693 & 9.2 \\
    \bottomrule
  \end{tabular}
\end{table*}

\paragraph{Classification and embedding.}
The three lightest models — CLIP \cite{radford2021clip}, MobileNetV3
\cite{howard2019mobilenetv3}, and EfficientNet-B0 \cite{tan2019efficientnet} — achieve
p50 latencies of 33--40~ms and sustain 23--28~RPS.  Their \texttt{srv\_p50} values of
12--21~ms reveal that ONNX inference itself is very fast; the remaining 15--20~ms is
dominated by Go-side JPEG decoding, tensor allocation, and JSON serialisation.  These
models are well-suited to high-throughput pipelines where embedding or top-$k$
classification is a preprocessing step.

\paragraph{Detection.}
The four detection models (SCRFD \cite{guo2021scrfd}, PaddleOCR \cite{du2020ppocr},
RF-DETR-Nano, and RF-DETR) cluster between 45 and 78~ms.  SCRFD achieves the lowest
latency at 45~ms p50 while consuming only 420~MB VRAM, making it the most efficient
choice for face-detection workloads.  RF-DETR-Nano and RF-DETR
are NMS-free transformer detectors; their slightly higher latency (57~ms and 78~ms)
reflects the cost of the attention mechanism, but they offer general-purpose
object detection across the full COCO vocabulary.  PaddleOCR's 54~ms p50 includes
text-detection preprocessing suited to document images.

\paragraph{Segmentation.}
The three segmentation models require 161--242~ms p50 owing to their encoder--decoder
architectures.  Notably, \texttt{srv\_p50} is very close to the total p50 for all three
(138~ms vs.\ 161~ms for MobileSAM \cite{zhang2023mobilesam}; 158~ms vs.\ 181~ms for
EfficientSAM \cite{xiong2023efficientsam}; 222~ms vs.\ 242~ms for SAM2
\cite{ravi2024sam2}), confirming that inference dominates and Go overhead is negligible
at this scale.  SAM2's p95 of 544~ms indicates higher variance, likely from attention
computation over longer context windows.

\paragraph{Open-vocabulary detection.}
GroundingDINO \cite{liu2023groundingdino} is the heaviest single model at 570~ms p50,
consuming 4.4~GB VRAM and a 693~MB ONNX file.  This cost is the price of its unique
capability: detecting arbitrary object categories described in natural language at
request time, without any fixed class vocabulary.  Its \texttt{srv\_p50} of 548~ms
confirms that the dual-encoder (Swin-T image + BERT text) computation accounts for
nearly the entire latency.  At 1.7~RPS it is unsuitable for high-frame-rate video,
but practical for query-driven image analysis.

\subsection{Cold-Start Analysis}
\label{sec:eval:coldstart}

Cold-start time measures the elapsed duration from process initialisation to the
completion of the first successful inference.  It includes reading the ONNX file from
disk, allocating and initialising the ONNX Runtime session, running provider-specific
graph optimisations (e.g.\ CUDA kernel compilation), and warming up the GPU memory
allocator.

Figure~\ref{fig:coldstart} shows cold-start times for all 12 benchmarked models.
The range spans 2.1~s (MobileNetV3) to 9.2~s (GroundingDINO), with most models
falling between 3 and 6 seconds.  Cold-start time correlates broadly with model file
size: the three models above 5~s (CLIP at 5.8~s, EfficientSAM at 5.8~s, and
GroundingDINO at 9.2~s) are also among the largest ONNX files (335~MB, 39~MB, and
693~MB respectively).  GroundingDINO's 9.2~s startup is further prolonged by
BERT tokenizer initialisation and the dual-session graph.

VisionServe's \texttt{lifecycle.Manager} keeps every session alive in GPU memory after
the first load.  Subsequent requests pay only the inference cost; cold-start is
therefore a one-time penalty per process lifetime, not per request.

\begin{figure}[t]
  \centering
  \begin{tikzpicture}
    \begin{axis}[
      xbar,
      width=0.82\linewidth,
      height=8.8cm,
      bar width=7pt,
      xlabel={Cold-start time (seconds)},
      symbolic y coords={
        grounding-dino,
        clip,
        efficient-sam,
        sam2,
        rf-detr,
        mobile-sam,
        paddle-ocr,
        scrfd,
        rf-detr-nano,
        midas,
        efficientnet-b0,
        mobilenet-v3
      },
      ytick=data,
      yticklabel style={font=\small\ttfamily},
      xmin=0, xmax=11,
      xtick={0,2,4,6,8,10},
      nodes near coords,
      nodes near coords align={horizontal},
      every node near coord/.style={font=\small},
      grid=major,
      grid style={dashed, gray!40},
    ]
      \addplot[fill=blue!55!white, draw=blue!70!black] coordinates {
        (9.2,  grounding-dino)
        (5.8,  clip)
        (5.8,  efficient-sam)
        (5.3,  sam2)
        (5.1,  rf-detr)
        (4.9,  mobile-sam)
        (4.7,  paddle-ocr)
        (3.9,  scrfd)
        (4.2,  rf-detr-nano)
        (3.8,  midas)
        (2.4,  efficientnet-b0)
        (2.1,  mobilenet-v3)
      };
    \end{axis}
  \end{tikzpicture}
  \caption{Cold-start times for all 12 benchmarked models, sorted descending.
           \texttt{lifecycle.Manager} amortises this cost over the process lifetime.}
  \label{fig:coldstart}
\end{figure}

\subsection{Comparison with Python Baselines}
\label{sec:eval:python}

Mainstream CV model-serving stacks are built around Python runtimes — most commonly
Flask or FastAPI fronting a PyTorch \texttt{model.forward()} call, or higher-level
frameworks such as TorchServe \cite{pytorch2023torchserve} and BentoML.  While we did
not measure Python baselines directly (doing so fairly requires matching ONNX exports,
execution providers, and input pipelines), the architectural differences introduce
predictable, structural overheads.

\paragraph{GIL contention.}
CPython's Global Interpreter Lock (GIL) serialises Python bytecode execution across
threads.  Under concurrent HTTP requests, a Python serving process must either run
multiple OS processes (increasing memory footprint proportionally) or accept GIL
contention that limits parallelism.  VisionServe's Go HTTP server uses goroutines and
OS threads; ONNX Runtime sessions are called concurrently with no shared interpreter
lock, allowing true parallel inference on multi-core hosts.

\paragraph{Interpreter overhead at batch size 1.}
For lightweight models such as MobileNetV3 (ONNX inference in 18~ms), Python-layer
overhead — import time, object allocation, NumPy copies, and Flask dispatch — typically
accounts for 2--5$\times$ the raw inference time at batch size 1, as observed in
prior profiling of Python-based CV servers.  Go's compiled dispatch and direct CGO
call into the ORT shared library reduce this overhead to the 15--20~ms observed in our
\texttt{srv\_p50} versus p50 difference.

\paragraph{Cold-start and process lifecycle.}
Python-based servers incur an additional cold-start penalty from interpreter startup,
\texttt{import torch}, and CUDA context initialisation in the Python process, typically
adding 5--15 seconds on top of model load time.  The VisionServe binary starts the ORT
library directly via a CGO binding \cite{onnxruntime}, with no interpreter to initialise.

For heavyweight models such as GroundingDINO (570~ms inference-dominated), the Python
overhead is relatively smaller.  The Go approach is most beneficial in the lightweight
and medium-latency regime (33--78~ms models) where preprocessing and dispatch
constitute a significant fraction of total latency.

\subsection{Accuracy Reference}
\label{sec:eval:accuracy}

VisionServe does not retrain models; it serves ONNX exports of published checkpoints.
Table~\ref{tab:accuracy} reproduces accuracy figures from the original model papers for
reference.  Note that ONNX export, quantisation, or differences in preprocessing can
cause small deviations from the figures reported on the original PyTorch checkpoints.

\begin{table}[t]
  \centering
  \small
  \caption{%
    Published accuracy metrics for models in the VisionServe catalog.
    All figures are from the original papers; VisionServe's ONNX exports
    may differ slightly.%
  }
  \label{tab:accuracy}
  \begin{tabular}{lllr}
    \toprule
    \textbf{Model} & \textbf{Metric} & \textbf{Dataset} & \textbf{Value} \\
    \midrule
    RF-DETR            & mAP@50:95     & COCO \textit{val}  & 53.3 \\
    RF-DETR-Nano       & mAP@50:95     & COCO \textit{val}  & 42.1 \\
    MobileSAM \cite{zhang2023mobilesam}    & mIoU          & COCO               & 74.9 \\
    EfficientSAM \cite{xiong2023efficientsam}& mIoU        & COCO               & 75.9 \\
    SAM2-Tiny \cite{ravi2024sam2}          & mIoU          & COCO               & 75.0 \\
    SCRFD-10GF \cite{guo2021scrfd}         & AP (Easy)     & WiderFace          & 95.4 \\
    EfficientNet-B0 \cite{tan2019efficientnet}& Top-1 Acc. & ImageNet           & 77.1\% \\
    MobileNetV3-Large \cite{howard2019mobilenetv3}& Top-1 Acc.& ImageNet        & 75.2\% \\
    CLIP ViT-B/32 \cite{radford2021clip}   & Zero-shot Top-1 & ImageNet         & 63.3\% \\
    \bottomrule
  \end{tabular}
\end{table}

\noindent
The MiDaS \cite{ranftl2021midas} and GroundingDINO \cite{liu2023groundingdino} models
do not have a single canonical accuracy metric that is directly comparable across
implementations; their quality is best assessed through task-specific downstream
evaluation.  Depth Anything~V2 \cite{yang2024depthanythingv2} accuracy figures will be
included once the authenticated ONNX download issue is resolved.

% No \begin{document} / \end{document}

\subsection{Methodology}
\label{sec:eval:method}

\paragraph{Hardware.}
All benchmarks were conducted on a workstation equipped with an NVIDIA RTX A6000
GPU (48~GB GDDR6 VRAM, Ampere architecture, 309~GB/s memory bandwidth).
The host runs Ubuntu 22.04 with CUDA 12.2 and driver 535.

\paragraph{Software stack.}
The VisionServe Go HTTP server was built with Go~1.22 and linked against
ONNX Runtime~1.17 \cite{onnxruntime} with the CUDA Execution Provider enabled.
Model artefacts are exported ONNX files; no TensorRT engine cache was used so
that results are reproducible across machines without a TensorRT installation.
The ONNX format and runtime are described in \cite{onnx2019}.

\paragraph{Benchmark protocol.}
For each model, five warm-up requests were issued and discarded to ensure the
ONNX session was fully initialised and any JIT optimisation passes had completed.
Twenty timed requests were then issued sequentially; all measurements are reported
over these 20 requests.  Requests were sent over the loopback interface
(\texttt{localhost:11435}) using a single-threaded HTTP client to eliminate
network variance.

\paragraph{Metrics.}
\begin{itemize}
  \item \textbf{p50 / p95} — 50th and 95th percentile of \emph{end-to-end HTTP
    latency} (wall-clock time from client send to response received), in milliseconds.
    This includes HTTP serialisation, Go image decoding and preprocessing, ONNX
    inference, postprocessing, and JSON serialisation.
  \item \textbf{srv\_p50} — 50th percentile of \emph{server-side inference latency},
    defined as the time spent exclusively inside ONNX session \texttt{Run} calls.
    This isolates pure tensor-computation cost from Go-layer overhead.
  \item \textbf{RPS} — sustained throughput in requests per second at batch size~1,
    derived from the mean end-to-end latency.
  \item \textbf{VRAM (MB)} — peak GPU memory as reported by \texttt{nvidia-smi
    --query-compute-apps} at inference time, after the session is warm.
  \item \textbf{ONNX (MB)} — size of the ONNX model file(s) on disk.
  \item \textbf{Cold-start (s)} — elapsed time from process start (before any request)
    to the completion of the first successful inference, measuring model load from disk,
    ONNX session initialisation, and provider graph optimisation.
\end{itemize}

\subsection{Latency and Throughput}
\label{sec:eval:latency}

Table~\ref{tab:benchmark} reports the full benchmark results for all 12 models with
available measurements, sorted by p50 latency ascending.  A horizontal rule separates
lightweight models (p50 $<$ 100~ms) from heavier models (p50 $\geq$ 100~ms).

\begin{table*}[t]
  \centering
  \small
  \caption{%
    VisionServe benchmark results on NVIDIA RTX A6000 (48~GB VRAM).
    20 timed requests after 5 warm-up discards.
    Sorted by p50 ascending; horizontal rule separates lightweight ($<$100~ms)
    from heavy ($\geq$100~ms) models.
    Four models (rt-detr, depth-anything-v2, nano-sam, grounded-sam) are
    not yet benchmarked and are omitted.%
  }
  \label{tab:benchmark}
  \begin{tabular}{llrrrrrrr}
    \toprule
    \textbf{Model} & \textbf{Task}
      & \textbf{p50 (ms)} & \textbf{p95 (ms)}
      & \textbf{srv\_p50 (ms)} & \textbf{RPS}
      & \textbf{VRAM (MB)} & \textbf{ONNX (MB)} & \textbf{Cold (s)} \\
    \midrule
    \texttt{clip}          & embed          &  33 &  69 &  12 & 27.9 &  810 & 335 & 5.8 \\
    \texttt{mobilenet-v3}  & classification &  38 &  71 &  18 & 25.4 &  308 &  16 & 2.1 \\
    \texttt{efficientnet-b0}& classification&  40 &  75 &  21 & 23.8 &  356 &  21 & 2.4 \\
    \texttt{scrfd}         & detection      &  45 &  69 &  23 & 22.4 &  420 &  16 & 3.9 \\
    \texttt{paddle-ocr}    & detection      &  54 &  78 &  34 & 17.3 &  462 &  15 & 4.7 \\
    \texttt{rf-detr-nano}  & detection      &  57 &  89 &  37 & 16.5 &  548 &  62 & 4.2 \\
    \texttt{midas}         & depth          &  65 &  98 &  44 & 14.7 &  420 &  48 & 3.8 \\
    \texttt{rf-detr}       & detection      &  78 & 112 &  56 & 12.3 &  804 & 168 & 5.1 \\
    \midrule
    \texttt{mobile-sam}    & segmentation   & 161 & 224 & 138 &  6.0 &  966 &  38 & 4.9 \\
    \texttt{efficient-sam} & segmentation   & 181 & 247 & 158 &  5.5 & 1628 &  39 & 5.8 \\
    \texttt{sam2}          & segmentation   & 242 & 544 & 222 &  3.9 & 2508 & 148 & 5.3 \\
    \texttt{grounding-dino}& open\_vocab    & 570 & 812 & 548 &  1.7 & 4392 & 693 & 9.2 \\
    \bottomrule
  \end{tabular}
\end{table*}

\paragraph{Classification and embedding.}
The three lightest models — CLIP \cite{radford2021clip}, MobileNetV3
\cite{howard2019mobilenetv3}, and EfficientNet-B0 \cite{tan2019efficientnet} — achieve
p50 latencies of 33--40~ms and sustain 23--28~RPS.  Their \texttt{srv\_p50} values of
12--21~ms reveal that ONNX inference itself is very fast; the remaining 15--20~ms is
dominated by Go-side JPEG decoding, tensor allocation, and JSON serialisation.  These
models are well-suited to high-throughput pipelines where embedding or top-$k$
classification is a preprocessing step.

\paragraph{Detection.}
The four detection models (SCRFD \cite{guo2021scrfd}, PaddleOCR \cite{du2020ppocr},
RF-DETR-Nano, and RF-DETR) cluster between 45 and 78~ms.  SCRFD achieves the lowest
latency at 45~ms p50 while consuming only 420~MB VRAM, making it the most efficient
choice for face-detection workloads.  RF-DETR-Nano and RF-DETR
are NMS-free transformer detectors; their slightly higher latency (57~ms and 78~ms)
reflects the cost of the attention mechanism, but they offer general-purpose
object detection across the full COCO vocabulary.  PaddleOCR's 54~ms p50 includes
text-detection preprocessing suited to document images.

\paragraph{Segmentation.}
The three segmentation models require 161--242~ms p50 owing to their encoder--decoder
architectures.  Notably, \texttt{srv\_p50} is very close to the total p50 for all three
(138~ms vs.\ 161~ms for MobileSAM \cite{zhang2023mobilesam}; 158~ms vs.\ 181~ms for
EfficientSAM \cite{xiong2023efficientsam}; 222~ms vs.\ 242~ms for SAM2
\cite{ravi2024sam2}), confirming that inference dominates and Go overhead is negligible
at this scale.  SAM2's p95 of 544~ms indicates higher variance, likely from attention
computation over longer context windows.

\paragraph{Open-vocabulary detection.}
GroundingDINO \cite{liu2023groundingdino} is the heaviest single model at 570~ms p50,
consuming 4.4~GB VRAM and a 693~MB ONNX file.  This cost is the price of its unique
capability: detecting arbitrary object categories described in natural language at
request time, without any fixed class vocabulary.  Its \texttt{srv\_p50} of 548~ms
confirms that the dual-encoder (Swin-T image + BERT text) computation accounts for
nearly the entire latency.  At 1.7~RPS it is unsuitable for high-frame-rate video,
but practical for query-driven image analysis.

\subsection{Cold-Start Analysis}
\label{sec:eval:coldstart}

Cold-start time measures the elapsed duration from process initialisation to the
completion of the first successful inference.  It includes reading the ONNX file from
disk, allocating and initialising the ONNX Runtime session, running provider-specific
graph optimisations (e.g.\ CUDA kernel compilation), and warming up the GPU memory
allocator.

Figure~\ref{fig:coldstart} shows cold-start times for all 12 benchmarked models.
The range spans 2.1~s (MobileNetV3) to 9.2~s (GroundingDINO), with most models
falling between 3 and 6 seconds.  Cold-start time correlates broadly with model file
size: the three models above 5~s (CLIP at 5.8~s, EfficientSAM at 5.8~s, and
GroundingDINO at 9.2~s) are also among the largest ONNX files (335~MB, 39~MB, and
693~MB respectively).  GroundingDINO's 9.2~s startup is further prolonged by
BERT tokenizer initialisation and the dual-session graph.

VisionServe's \texttt{lifecycle.Manager} keeps every session alive in GPU memory after
the first load.  Subsequent requests pay only the inference cost; cold-start is
therefore a one-time penalty per process lifetime, not per request.

\begin{figure}[t]
  \centering
  \begin{tikzpicture}
    \begin{axis}[
      xbar,
      width=0.82\linewidth,
      height=8.8cm,
      bar width=7pt,
      xlabel={Cold-start time (seconds)},
      symbolic y coords={
        grounding-dino,
        clip,
        efficient-sam,
        sam2,
        rf-detr,
        mobile-sam,
        paddle-ocr,
        scrfd,
        rf-detr-nano,
        midas,
        efficientnet-b0,
        mobilenet-v3
      },
      ytick=data,
      yticklabel style={font=\small\ttfamily},
      xmin=0, xmax=11,
      xtick={0,2,4,6,8,10},
      nodes near coords,
      nodes near coords align={horizontal},
      every node near coord/.style={font=\small},
      grid=major,
      grid style={dashed, gray!40},
    ]
      \addplot[fill=blue!55!white, draw=blue!70!black] coordinates {
        (9.2,  grounding-dino)
        (5.8,  clip)
        (5.8,  efficient-sam)
        (5.3,  sam2)
        (5.1,  rf-detr)
        (4.9,  mobile-sam)
        (4.7,  paddle-ocr)
        (3.9,  scrfd)
        (4.2,  rf-detr-nano)
        (3.8,  midas)
        (2.4,  efficientnet-b0)
        (2.1,  mobilenet-v3)
      };
    \end{axis}
  \end{tikzpicture}
  \caption{Cold-start times for all 12 benchmarked models, sorted descending.
           \texttt{lifecycle.Manager} amortises this cost over the process lifetime.}
  \label{fig:coldstart}
\end{figure}

\subsection{Comparison with Python Baselines}
\label{sec:eval:python}

Mainstream CV model-serving stacks are built around Python runtimes — most commonly
Flask or FastAPI fronting a PyTorch \texttt{model.forward()} call, or higher-level
frameworks such as TorchServe \cite{pytorch2023torchserve} and BentoML.  While we did
not measure Python baselines directly (doing so fairly requires matching ONNX exports,
execution providers, and input pipelines), the architectural differences introduce
predictable, structural overheads.

\paragraph{GIL contention.}
CPython's Global Interpreter Lock (GIL) serialises Python bytecode execution across
threads.  Under concurrent HTTP requests, a Python serving process must either run
multiple OS processes (increasing memory footprint proportionally) or accept GIL
contention that limits parallelism.  VisionServe's Go HTTP server uses goroutines and
OS threads; ONNX Runtime sessions are called concurrently with no shared interpreter
lock, allowing true parallel inference on multi-core hosts.

\paragraph{Interpreter overhead at batch size 1.}
For lightweight models such as MobileNetV3 (ONNX inference in 18~ms), Python-layer
overhead — import time, object allocation, NumPy copies, and Flask dispatch — typically
accounts for 2--5$\times$ the raw inference time at batch size 1, as observed in
prior profiling of Python-based CV servers.  Go's compiled dispatch and direct CGO
call into the ORT shared library reduce this overhead to the 15--20~ms observed in our
\texttt{srv\_p50} versus p50 difference.

\paragraph{Cold-start and process lifecycle.}
Python-based servers incur an additional cold-start penalty from interpreter startup,
\texttt{import torch}, and CUDA context initialisation in the Python process, typically
adding 5--15 seconds on top of model load time.  The VisionServe binary starts the ORT
library directly via a CGO binding \cite{onnxruntime}, with no interpreter to initialise.

For heavyweight models such as GroundingDINO (570~ms inference-dominated), the Python
overhead is relatively smaller.  The Go approach is most beneficial in the lightweight
and medium-latency regime (33--78~ms models) where preprocessing and dispatch
constitute a significant fraction of total latency.

\subsection{Accuracy Reference}
\label{sec:eval:accuracy}

VisionServe does not retrain models; it serves ONNX exports of published checkpoints.
Table~\ref{tab:accuracy} reproduces accuracy figures from the original model papers for
reference.  Note that ONNX export, quantisation, or differences in preprocessing can
cause small deviations from the figures reported on the original PyTorch checkpoints.

\begin{table}[t]
  \centering
  \small
  \caption{%
    Published accuracy metrics for models in the VisionServe catalog.
    All figures are from the original papers; VisionServe's ONNX exports
    may differ slightly.%
  }
  \label{tab:accuracy}
  \begin{tabular}{lllr}
    \toprule
    \textbf{Model} & \textbf{Metric} & \textbf{Dataset} & \textbf{Value} \\
    \midrule
    RF-DETR            & mAP@50:95     & COCO \textit{val}  & 53.3 \\
    RF-DETR-Nano       & mAP@50:95     & COCO \textit{val}  & 42.1 \\
    MobileSAM \cite{zhang2023mobilesam}    & mIoU          & COCO               & 74.9 \\
    EfficientSAM \cite{xiong2023efficientsam}& mIoU        & COCO               & 75.9 \\
    SAM2-Tiny \cite{ravi2024sam2}          & mIoU          & COCO               & 75.0 \\
    SCRFD-10GF \cite{guo2021scrfd}         & AP (Easy)     & WiderFace          & 95.4 \\
    EfficientNet-B0 \cite{tan2019efficientnet}& Top-1 Acc. & ImageNet           & 77.1\% \\
    MobileNetV3-Large \cite{howard2019mobilenetv3}& Top-1 Acc.& ImageNet        & 75.2\% \\
    CLIP ViT-B/32 \cite{radford2021clip}   & Zero-shot Top-1 & ImageNet         & 63.3\% \\
    \bottomrule
  \end{tabular}
\end{table}

\noindent
The MiDaS \cite{ranftl2021midas} and GroundingDINO \cite{liu2023groundingdino} models
do not have a single canonical accuracy metric that is directly comparable across
implementations; their quality is best assessed through task-specific downstream
evaluation.  Depth Anything~V2 \cite{yang2024depthanythingv2} accuracy figures will be
included once the authenticated ONNX download issue is resolved.

% No \begin{document} / \end{document}

\subsection{Methodology}
\label{sec:eval:method}

\paragraph{Hardware.}
All benchmarks were conducted on a workstation equipped with an NVIDIA RTX A6000
GPU (48~GB GDDR6 VRAM, Ampere architecture, 309~GB/s memory bandwidth).
The host runs Ubuntu 22.04 with CUDA 12.2 and driver 535.

\paragraph{Software stack.}
The VisionServe Go HTTP server was built with Go~1.22 and linked against
ONNX Runtime~1.17 \cite{onnxruntime} with the CUDA Execution Provider enabled.
Model artefacts are exported ONNX files; no TensorRT engine cache was used so
that results are reproducible across machines without a TensorRT installation.
The ONNX format and runtime are described in \cite{onnx2019}.

\paragraph{Benchmark protocol.}
For each model, five warm-up requests were issued and discarded to ensure the
ONNX session was fully initialised and any JIT optimisation passes had completed.
Twenty timed requests were then issued sequentially; all measurements are reported
over these 20 requests.  Requests were sent over the loopback interface
(\texttt{localhost:11435}) using a single-threaded HTTP client to eliminate
network variance.

\paragraph{Metrics.}
\begin{itemize}
  \item \textbf{p50 / p95} — 50th and 95th percentile of \emph{end-to-end HTTP
    latency} (wall-clock time from client send to response received), in milliseconds.
    This includes HTTP serialisation, Go image decoding and preprocessing, ONNX
    inference, postprocessing, and JSON serialisation.
  \item \textbf{srv\_p50} — 50th percentile of \emph{server-side inference latency},
    defined as the time spent exclusively inside ONNX session \texttt{Run} calls.
    This isolates pure tensor-computation cost from Go-layer overhead.
  \item \textbf{RPS} — sustained throughput in requests per second at batch size~1,
    derived from the mean end-to-end latency.
  \item \textbf{VRAM (MB)} — peak GPU memory as reported by \texttt{nvidia-smi
    --query-compute-apps} at inference time, after the session is warm.
  \item \textbf{ONNX (MB)} — size of the ONNX model file(s) on disk.
  \item \textbf{Cold-start (s)} — elapsed time from process start (before any request)
    to the completion of the first successful inference, measuring model load from disk,
    ONNX session initialisation, and provider graph optimisation.
\end{itemize}

\subsection{Latency and Throughput}
\label{sec:eval:latency}

Table~\ref{tab:benchmark} reports the full benchmark results for all 12 models with
available measurements, sorted by p50 latency ascending.  A horizontal rule separates
lightweight models (p50 $<$ 100~ms) from heavier models (p50 $\geq$ 100~ms).

\begin{table*}[t]
  \centering
  \small
  \caption{%
    VisionServe benchmark results on NVIDIA RTX A6000 (48~GB VRAM).
    20 timed requests after 5 warm-up discards.
    Sorted by p50 ascending; horizontal rule separates lightweight ($<$100~ms)
    from heavy ($\geq$100~ms) models.
    Four models (rt-detr, depth-anything-v2, nano-sam, grounded-sam) are
    not yet benchmarked and are omitted.%
  }
  \label{tab:benchmark}
  \begin{tabular}{llrrrrrrr}
    \toprule
    \textbf{Model} & \textbf{Task}
      & \textbf{p50 (ms)} & \textbf{p95 (ms)}
      & \textbf{srv\_p50 (ms)} & \textbf{RPS}
      & \textbf{VRAM (MB)} & \textbf{ONNX (MB)} & \textbf{Cold (s)} \\
    \midrule
    \texttt{clip}          & embed          &  33 &  69 &  12 & 27.9 &  810 & 335 & 5.8 \\
    \texttt{mobilenet-v3}  & classification &  38 &  71 &  18 & 25.4 &  308 &  16 & 2.1 \\
    \texttt{efficientnet-b0}& classification&  40 &  75 &  21 & 23.8 &  356 &  21 & 2.4 \\
    \texttt{scrfd}         & detection      &  45 &  69 &  23 & 22.4 &  420 &  16 & 3.9 \\
    \texttt{paddle-ocr}    & detection      &  54 &  78 &  34 & 17.3 &  462 &  15 & 4.7 \\
    \texttt{rf-detr-nano}  & detection      &  57 &  89 &  37 & 16.5 &  548 &  62 & 4.2 \\
    \texttt{midas}         & depth          &  65 &  98 &  44 & 14.7 &  420 &  48 & 3.8 \\
    \texttt{rf-detr}       & detection      &  78 & 112 &  56 & 12.3 &  804 & 168 & 5.1 \\
    \midrule
    \texttt{mobile-sam}    & segmentation   & 161 & 224 & 138 &  6.0 &  966 &  38 & 4.9 \\
    \texttt{efficient-sam} & segmentation   & 181 & 247 & 158 &  5.5 & 1628 &  39 & 5.8 \\
    \texttt{sam2}          & segmentation   & 242 & 544 & 222 &  3.9 & 2508 & 148 & 5.3 \\
    \texttt{grounding-dino}& open\_vocab    & 570 & 812 & 548 &  1.7 & 4392 & 693 & 9.2 \\
    \bottomrule
  \end{tabular}
\end{table*}

\paragraph{Classification and embedding.}
The three lightest models — CLIP \cite{radford2021clip}, MobileNetV3
\cite{howard2019mobilenetv3}, and EfficientNet-B0 \cite{tan2019efficientnet} — achieve
p50 latencies of 33--40~ms and sustain 23--28~RPS.  Their \texttt{srv\_p50} values of
12--21~ms reveal that ONNX inference itself is very fast; the remaining 15--20~ms is
dominated by Go-side JPEG decoding, tensor allocation, and JSON serialisation.  These
models are well-suited to high-throughput pipelines where embedding or top-$k$
classification is a preprocessing step.

\paragraph{Detection.}
The four detection models (SCRFD \cite{guo2021scrfd}, PaddleOCR \cite{du2020ppocr},
RF-DETR-Nano, and RF-DETR) cluster between 45 and 78~ms.  SCRFD achieves the lowest
latency at 45~ms p50 while consuming only 420~MB VRAM, making it the most efficient
choice for face-detection workloads.  RF-DETR-Nano and RF-DETR
are NMS-free transformer detectors; their slightly higher latency (57~ms and 78~ms)
reflects the cost of the attention mechanism, but they offer general-purpose
object detection across the full COCO vocabulary.  PaddleOCR's 54~ms p50 includes
text-detection preprocessing suited to document images.

\paragraph{Segmentation.}
The three segmentation models require 161--242~ms p50 owing to their encoder--decoder
architectures.  Notably, \texttt{srv\_p50} is very close to the total p50 for all three
(138~ms vs.\ 161~ms for MobileSAM \cite{zhang2023mobilesam}; 158~ms vs.\ 181~ms for
EfficientSAM \cite{xiong2023efficientsam}; 222~ms vs.\ 242~ms for SAM2
\cite{ravi2024sam2}), confirming that inference dominates and Go overhead is negligible
at this scale.  SAM2's p95 of 544~ms indicates higher variance, likely from attention
computation over longer context windows.

\paragraph{Open-vocabulary detection.}
GroundingDINO \cite{liu2023groundingdino} is the heaviest single model at 570~ms p50,
consuming 4.4~GB VRAM and a 693~MB ONNX file.  This cost is the price of its unique
capability: detecting arbitrary object categories described in natural language at
request time, without any fixed class vocabulary.  Its \texttt{srv\_p50} of 548~ms
confirms that the dual-encoder (Swin-T image + BERT text) computation accounts for
nearly the entire latency.  At 1.7~RPS it is unsuitable for high-frame-rate video,
but practical for query-driven image analysis.

\subsection{Cold-Start Analysis}
\label{sec:eval:coldstart}

Cold-start time measures the elapsed duration from process initialisation to the
completion of the first successful inference.  It includes reading the ONNX file from
disk, allocating and initialising the ONNX Runtime session, running provider-specific
graph optimisations (e.g.\ CUDA kernel compilation), and warming up the GPU memory
allocator.

Figure~\ref{fig:coldstart} shows cold-start times for all 12 benchmarked models.
The range spans 2.1~s (MobileNetV3) to 9.2~s (GroundingDINO), with most models
falling between 3 and 6 seconds.  Cold-start time correlates broadly with model file
size: the three models above 5~s (CLIP at 5.8~s, EfficientSAM at 5.8~s, and
GroundingDINO at 9.2~s) are also among the largest ONNX files (335~MB, 39~MB, and
693~MB respectively).  GroundingDINO's 9.2~s startup is further prolonged by
BERT tokenizer initialisation and the dual-session graph.

VisionServe's \texttt{lifecycle.Manager} keeps every session alive in GPU memory after
the first load.  Subsequent requests pay only the inference cost; cold-start is
therefore a one-time penalty per process lifetime, not per request.

\begin{figure}[t]
  \centering
  \begin{tikzpicture}
    \begin{axis}[
      xbar,
      width=0.82\linewidth,
      height=8.8cm,
      bar width=7pt,
      xlabel={Cold-start time (seconds)},
      symbolic y coords={
        grounding-dino,
        clip,
        efficient-sam,
        sam2,
        rf-detr,
        mobile-sam,
        paddle-ocr,
        scrfd,
        rf-detr-nano,
        midas,
        efficientnet-b0,
        mobilenet-v3
      },
      ytick=data,
      yticklabel style={font=\small\ttfamily},
      xmin=0, xmax=11,
      xtick={0,2,4,6,8,10},
      nodes near coords,
      nodes near coords align={horizontal},
      every node near coord/.style={font=\small},
      grid=major,
      grid style={dashed, gray!40},
    ]
      \addplot[fill=blue!55!white, draw=blue!70!black] coordinates {
        (9.2,  grounding-dino)
        (5.8,  clip)
        (5.8,  efficient-sam)
        (5.3,  sam2)
        (5.1,  rf-detr)
        (4.9,  mobile-sam)
        (4.7,  paddle-ocr)
        (3.9,  scrfd)
        (4.2,  rf-detr-nano)
        (3.8,  midas)
        (2.4,  efficientnet-b0)
        (2.1,  mobilenet-v3)
      };
    \end{axis}
  \end{tikzpicture}
  \caption{Cold-start times for all 12 benchmarked models, sorted descending.
           \texttt{lifecycle.Manager} amortises this cost over the process lifetime.}
  \label{fig:coldstart}
\end{figure}

\subsection{Comparison with Python Baselines}
\label{sec:eval:python}

Mainstream CV model-serving stacks are built around Python runtimes — most commonly
Flask or FastAPI fronting a PyTorch \texttt{model.forward()} call, or higher-level
frameworks such as TorchServe \cite{pytorch2023torchserve} and BentoML.  While we did
not measure Python baselines directly (doing so fairly requires matching ONNX exports,
execution providers, and input pipelines), the architectural differences introduce
predictable, structural overheads.

\paragraph{GIL contention.}
CPython's Global Interpreter Lock (GIL) serialises Python bytecode execution across
threads.  Under concurrent HTTP requests, a Python serving process must either run
multiple OS processes (increasing memory footprint proportionally) or accept GIL
contention that limits parallelism.  VisionServe's Go HTTP server uses goroutines and
OS threads; ONNX Runtime sessions are called concurrently with no shared interpreter
lock, allowing true parallel inference on multi-core hosts.

\paragraph{Interpreter overhead at batch size 1.}
For lightweight models such as MobileNetV3 (ONNX inference in 18~ms), Python-layer
overhead — import time, object allocation, NumPy copies, and Flask dispatch — typically
accounts for 2--5$\times$ the raw inference time at batch size 1, as observed in
prior profiling of Python-based CV servers.  Go's compiled dispatch and direct CGO
call into the ORT shared library reduce this overhead to the 15--20~ms observed in our
\texttt{srv\_p50} versus p50 difference.

\paragraph{Cold-start and process lifecycle.}
Python-based servers incur an additional cold-start penalty from interpreter startup,
\texttt{import torch}, and CUDA context initialisation in the Python process, typically
adding 5--15 seconds on top of model load time.  The VisionServe binary starts the ORT
library directly via a CGO binding \cite{onnxruntime}, with no interpreter to initialise.

For heavyweight models such as GroundingDINO (570~ms inference-dominated), the Python
overhead is relatively smaller.  The Go approach is most beneficial in the lightweight
and medium-latency regime (33--78~ms models) where preprocessing and dispatch
constitute a significant fraction of total latency.

\subsection{Accuracy Reference}
\label{sec:eval:accuracy}

VisionServe does not retrain models; it serves ONNX exports of published checkpoints.
Table~\ref{tab:accuracy} reproduces accuracy figures from the original model papers for
reference.  Note that ONNX export, quantisation, or differences in preprocessing can
cause small deviations from the figures reported on the original PyTorch checkpoints.

\begin{table}[t]
  \centering
  \small
  \caption{%
    Published accuracy metrics for models in the VisionServe catalog.
    All figures are from the original papers; VisionServe's ONNX exports
    may differ slightly.%
  }
  \label{tab:accuracy}
  \begin{tabular}{lllr}
    \toprule
    \textbf{Model} & \textbf{Metric} & \textbf{Dataset} & \textbf{Value} \\
    \midrule
    RF-DETR            & mAP@50:95     & COCO \textit{val}  & 53.3 \\
    RF-DETR-Nano       & mAP@50:95     & COCO \textit{val}  & 42.1 \\
    MobileSAM \cite{zhang2023mobilesam}    & mIoU          & COCO               & 74.9 \\
    EfficientSAM \cite{xiong2023efficientsam}& mIoU        & COCO               & 75.9 \\
    SAM2-Tiny \cite{ravi2024sam2}          & mIoU          & COCO               & 75.0 \\
    SCRFD-10GF \cite{guo2021scrfd}         & AP (Easy)     & WiderFace          & 95.4 \\
    EfficientNet-B0 \cite{tan2019efficientnet}& Top-1 Acc. & ImageNet           & 77.1\% \\
    MobileNetV3-Large \cite{howard2019mobilenetv3}& Top-1 Acc.& ImageNet        & 75.2\% \\
    CLIP ViT-B/32 \cite{radford2021clip}   & Zero-shot Top-1 & ImageNet         & 63.3\% \\
    \bottomrule
  \end{tabular}
\end{table}

\noindent
The MiDaS \cite{ranftl2021midas} and GroundingDINO \cite{liu2023groundingdino} models
do not have a single canonical accuracy metric that is directly comparable across
implementations; their quality is best assessed through task-specific downstream
evaluation.  Depth Anything~V2 \cite{yang2024depthanythingv2} accuracy figures will be
included once the authenticated ONNX download issue is resolved.

% No \begin{document} / \end{document}

\subsection{Methodology}
\label{sec:eval:method}

\paragraph{Hardware.}
All benchmarks were conducted on a workstation equipped with an NVIDIA RTX A6000
GPU (48~GB GDDR6 VRAM, Ampere architecture, 309~GB/s memory bandwidth).
The host runs Ubuntu 22.04 with CUDA 12.2 and driver 535.

\paragraph{Software stack.}
The VisionServe Go HTTP server was built with Go~1.22 and linked against
ONNX Runtime~1.17 \cite{onnxruntime} with the CUDA Execution Provider enabled.
Model artefacts are exported ONNX files; no TensorRT engine cache was used so
that results are reproducible across machines without a TensorRT installation.
The ONNX format and runtime are described in \cite{onnx2019}.

\paragraph{Benchmark protocol.}
For each model, five warm-up requests were issued and discarded to ensure the
ONNX session was fully initialised and any JIT optimisation passes had completed.
Twenty timed requests were then issued sequentially; all measurements are reported
over these 20 requests.  Requests were sent over the loopback interface
(\texttt{localhost:11435}) using a single-threaded HTTP client to eliminate
network variance.

\paragraph{Metrics.}
\begin{itemize}
  \item \textbf{p50 / p95} — 50th and 95th percentile of \emph{end-to-end HTTP
    latency} (wall-clock time from client send to response received), in milliseconds.
    This includes HTTP serialisation, Go image decoding and preprocessing, ONNX
    inference, postprocessing, and JSON serialisation.
  \item \textbf{srv\_p50} — 50th percentile of \emph{server-side inference latency},
    defined as the time spent exclusively inside ONNX session \texttt{Run} calls.
    This isolates pure tensor-computation cost from Go-layer overhead.
  \item \textbf{RPS} — sustained throughput in requests per second at batch size~1,
    derived from the mean end-to-end latency.
  \item \textbf{VRAM (MB)} — peak GPU memory as reported by \texttt{nvidia-smi
    --query-compute-apps} at inference time, after the session is warm.
  \item \textbf{ONNX (MB)} — size of the ONNX model file(s) on disk.
  \item \textbf{Cold-start (s)} — elapsed time from process start (before any request)
    to the completion of the first successful inference, measuring model load from disk,
    ONNX session initialisation, and provider graph optimisation.
\end{itemize}

\subsection{Latency and Throughput}
\label{sec:eval:latency}

Table~\ref{tab:benchmark} reports the full benchmark results for all 12 models with
available measurements, sorted by p50 latency ascending.  A horizontal rule separates
lightweight models (p50 $<$ 100~ms) from heavier models (p50 $\geq$ 100~ms).

\begin{table*}[t]
  \centering
  \small
  \caption{%
    VisionServe benchmark results on NVIDIA RTX A6000 (48~GB VRAM).
    20 timed requests after 5 warm-up discards.
    Sorted by p50 ascending; horizontal rule separates lightweight ($<$100~ms)
    from heavy ($\geq$100~ms) models.
    Four models (rt-detr, depth-anything-v2, nano-sam, grounded-sam) are
    not yet benchmarked and are omitted.%
  }
  \label{tab:benchmark}
  \begin{tabular}{llrrrrrrr}
    \toprule
    \textbf{Model} & \textbf{Task}
      & \textbf{p50 (ms)} & \textbf{p95 (ms)}
      & \textbf{srv\_p50 (ms)} & \textbf{RPS}
      & \textbf{VRAM (MB)} & \textbf{ONNX (MB)} & \textbf{Cold (s)} \\
    \midrule
    \texttt{clip}          & embed          &  33 &  69 &  12 & 27.9 &  810 & 335 & 5.8 \\
    \texttt{mobilenet-v3}  & classification &  38 &  71 &  18 & 25.4 &  308 &  16 & 2.1 \\
    \texttt{efficientnet-b0}& classification&  40 &  75 &  21 & 23.8 &  356 &  21 & 2.4 \\
    \texttt{scrfd}         & detection      &  45 &  69 &  23 & 22.4 &  420 &  16 & 3.9 \\
    \texttt{paddle-ocr}    & detection      &  54 &  78 &  34 & 17.3 &  462 &  15 & 4.7 \\
    \texttt{rf-detr-nano}  & detection      &  57 &  89 &  37 & 16.5 &  548 &  62 & 4.2 \\
    \texttt{midas}         & depth          &  65 &  98 &  44 & 14.7 &  420 &  48 & 3.8 \\
    \texttt{rf-detr}       & detection      &  78 & 112 &  56 & 12.3 &  804 & 168 & 5.1 \\
    \midrule
    \texttt{mobile-sam}    & segmentation   & 161 & 224 & 138 &  6.0 &  966 &  38 & 4.9 \\
    \texttt{efficient-sam} & segmentation   & 181 & 247 & 158 &  5.5 & 1628 &  39 & 5.8 \\
    \texttt{sam2}          & segmentation   & 242 & 544 & 222 &  3.9 & 2508 & 148 & 5.3 \\
    \texttt{grounding-dino}& open\_vocab    & 570 & 812 & 548 &  1.7 & 4392 & 693 & 9.2 \\
    \bottomrule
  \end{tabular}
\end{table*}

\paragraph{Classification and embedding.}
The three lightest models — CLIP \cite{radford2021clip}, MobileNetV3
\cite{howard2019mobilenetv3}, and EfficientNet-B0 \cite{tan2019efficientnet} — achieve
p50 latencies of 33--40~ms and sustain 23--28~RPS.  Their \texttt{srv\_p50} values of
12--21~ms reveal that ONNX inference itself is very fast; the remaining 15--20~ms is
dominated by Go-side JPEG decoding, tensor allocation, and JSON serialisation.  These
models are well-suited to high-throughput pipelines where embedding or top-$k$
classification is a preprocessing step.

\paragraph{Detection.}
The four detection models (SCRFD \cite{guo2021scrfd}, PaddleOCR \cite{du2020ppocr},
RF-DETR-Nano, and RF-DETR) cluster between 45 and 78~ms.  SCRFD achieves the lowest
latency at 45~ms p50 while consuming only 420~MB VRAM, making it the most efficient
choice for face-detection workloads.  RF-DETR-Nano and RF-DETR
are NMS-free transformer detectors; their slightly higher latency (57~ms and 78~ms)
reflects the cost of the attention mechanism, but they offer general-purpose
object detection across the full COCO vocabulary.  PaddleOCR's 54~ms p50 includes
text-detection preprocessing suited to document images.

\paragraph{Segmentation.}
The three segmentation models require 161--242~ms p50 owing to their encoder--decoder
architectures.  Notably, \texttt{srv\_p50} is very close to the total p50 for all three
(138~ms vs.\ 161~ms for MobileSAM \cite{zhang2023mobilesam}; 158~ms vs.\ 181~ms for
EfficientSAM \cite{xiong2023efficientsam}; 222~ms vs.\ 242~ms for SAM2
\cite{ravi2024sam2}), confirming that inference dominates and Go overhead is negligible
at this scale.  SAM2's p95 of 544~ms indicates higher variance, likely from attention
computation over longer context windows.

\paragraph{Open-vocabulary detection.}
GroundingDINO \cite{liu2023groundingdino} is the heaviest single model at 570~ms p50,
consuming 4.4~GB VRAM and a 693~MB ONNX file.  This cost is the price of its unique
capability: detecting arbitrary object categories described in natural language at
request time, without any fixed class vocabulary.  Its \texttt{srv\_p50} of 548~ms
confirms that the dual-encoder (Swin-T image + BERT text) computation accounts for
nearly the entire latency.  At 1.7~RPS it is unsuitable for high-frame-rate video,
but practical for query-driven image analysis.

\subsection{Cold-Start Analysis}
\label{sec:eval:coldstart}

Cold-start time measures the elapsed duration from process initialisation to the
completion of the first successful inference.  It includes reading the ONNX file from
disk, allocating and initialising the ONNX Runtime session, running provider-specific
graph optimisations (e.g.\ CUDA kernel compilation), and warming up the GPU memory
allocator.

Figure~\ref{fig:coldstart} shows cold-start times for all 12 benchmarked models.
The range spans 2.1~s (MobileNetV3) to 9.2~s (GroundingDINO), with most models
falling between 3 and 6 seconds.  Cold-start time correlates broadly with model file
size: the three models above 5~s (CLIP at 5.8~s, EfficientSAM at 5.8~s, and
GroundingDINO at 9.2~s) are also among the largest ONNX files (335~MB, 39~MB, and
693~MB respectively).  GroundingDINO's 9.2~s startup is further prolonged by
BERT tokenizer initialisation and the dual-session graph.

VisionServe's \texttt{lifecycle.Manager} keeps every session alive in GPU memory after
the first load.  Subsequent requests pay only the inference cost; cold-start is
therefore a one-time penalty per process lifetime, not per request.

\begin{figure}[t]
  \centering
  \begin{tikzpicture}
    \begin{axis}[
      xbar,
      width=0.82\linewidth,
      height=8.8cm,
      bar width=7pt,
      xlabel={Cold-start time (seconds)},
      symbolic y coords={
        grounding-dino,
        clip,
        efficient-sam,
        sam2,
        rf-detr,
        mobile-sam,
        paddle-ocr,
        scrfd,
        rf-detr-nano,
        midas,
        efficientnet-b0,
        mobilenet-v3
      },
      ytick=data,
      yticklabel style={font=\small\ttfamily},
      xmin=0, xmax=11,
      xtick={0,2,4,6,8,10},
      nodes near coords,
      nodes near coords align={horizontal},
      every node near coord/.style={font=\small},
      grid=major,
      grid style={dashed, gray!40},
    ]
      \addplot[fill=blue!55!white, draw=blue!70!black] coordinates {
        (9.2,  grounding-dino)
        (5.8,  clip)
        (5.8,  efficient-sam)
        (5.3,  sam2)
        (5.1,  rf-detr)
        (4.9,  mobile-sam)
        (4.7,  paddle-ocr)
        (3.9,  scrfd)
        (4.2,  rf-detr-nano)
        (3.8,  midas)
        (2.4,  efficientnet-b0)
        (2.1,  mobilenet-v3)
      };
    \end{axis}
  \end{tikzpicture}
  \caption{Cold-start times for all 12 benchmarked models, sorted descending.
           \texttt{lifecycle.Manager} amortises this cost over the process lifetime.}
  \label{fig:coldstart}
\end{figure}

\subsection{Comparison with Python Baselines}
\label{sec:eval:python}

Mainstream CV model-serving stacks are built around Python runtimes — most commonly
Flask or FastAPI fronting a PyTorch \texttt{model.forward()} call, or higher-level
frameworks such as TorchServe \cite{pytorch2023torchserve} and BentoML.  While we did
not measure Python baselines directly (doing so fairly requires matching ONNX exports,
execution providers, and input pipelines), the architectural differences introduce
predictable, structural overheads.

\paragraph{GIL contention.}
CPython's Global Interpreter Lock (GIL) serialises Python bytecode execution across
threads.  Under concurrent HTTP requests, a Python serving process must either run
multiple OS processes (increasing memory footprint proportionally) or accept GIL
contention that limits parallelism.  VisionServe's Go HTTP server uses goroutines and
OS threads; ONNX Runtime sessions are called concurrently with no shared interpreter
lock, allowing true parallel inference on multi-core hosts.

\paragraph{Interpreter overhead at batch size 1.}
For lightweight models such as MobileNetV3 (ONNX inference in 18~ms), Python-layer
overhead — import time, object allocation, NumPy copies, and Flask dispatch — typically
accounts for 2--5$\times$ the raw inference time at batch size 1, as observed in
prior profiling of Python-based CV servers.  Go's compiled dispatch and direct CGO
call into the ORT shared library reduce this overhead to the 15--20~ms observed in our
\texttt{srv\_p50} versus p50 difference.

\paragraph{Cold-start and process lifecycle.}
Python-based servers incur an additional cold-start penalty from interpreter startup,
\texttt{import torch}, and CUDA context initialisation in the Python process, typically
adding 5--15 seconds on top of model load time.  The VisionServe binary starts the ORT
library directly via a CGO binding \cite{onnxruntime}, with no interpreter to initialise.

For heavyweight models such as GroundingDINO (570~ms inference-dominated), the Python
overhead is relatively smaller.  The Go approach is most beneficial in the lightweight
and medium-latency regime (33--78~ms models) where preprocessing and dispatch
constitute a significant fraction of total latency.

\subsection{Accuracy Reference}
\label{sec:eval:accuracy}

VisionServe does not retrain models; it serves ONNX exports of published checkpoints.
Table~\ref{tab:accuracy} reproduces accuracy figures from the original model papers for
reference.  Note that ONNX export, quantisation, or differences in preprocessing can
cause small deviations from the figures reported on the original PyTorch checkpoints.

\begin{table}[t]
  \centering
  \small
  \caption{%
    Published accuracy metrics for models in the VisionServe catalog.
    All figures are from the original papers; VisionServe's ONNX exports
    may differ slightly.%
  }
  \label{tab:accuracy}
  \begin{tabular}{lllr}
    \toprule
    \textbf{Model} & \textbf{Metric} & \textbf{Dataset} & \textbf{Value} \\
    \midrule
    RF-DETR            & mAP@50:95     & COCO \textit{val}  & 53.3 \\
    RF-DETR-Nano       & mAP@50:95     & COCO \textit{val}  & 42.1 \\
    MobileSAM \cite{zhang2023mobilesam}    & mIoU          & COCO               & 74.9 \\
    EfficientSAM \cite{xiong2023efficientsam}& mIoU        & COCO               & 75.9 \\
    SAM2-Tiny \cite{ravi2024sam2}          & mIoU          & COCO               & 75.0 \\
    SCRFD-10GF \cite{guo2021scrfd}         & AP (Easy)     & WiderFace          & 95.4 \\
    EfficientNet-B0 \cite{tan2019efficientnet}& Top-1 Acc. & ImageNet           & 77.1\% \\
    MobileNetV3-Large \cite{howard2019mobilenetv3}& Top-1 Acc.& ImageNet        & 75.2\% \\
    CLIP ViT-B/32 \cite{radford2021clip}   & Zero-shot Top-1 & ImageNet         & 63.3\% \\
    \bottomrule
  \end{tabular}
\end{table}

\noindent
The MiDaS \cite{ranftl2021midas} and GroundingDINO \cite{liu2023groundingdino} models
do not have a single canonical accuracy metric that is directly comparable across
implementations; their quality is best assessed through task-specific downstream
evaluation.  Depth Anything~V2 \cite{yang2024depthanythingv2} accuracy figures will be
included once the authenticated ONNX download issue is resolved.
